\documentclass[11pt,a4paper]{article}
\usepackage[margin=1in]{geometry}
\usepackage{amsmath,booktabs,hyperref,float,xcolor,graphicx,multirow}
\definecolor{good}{rgb}{0.0,0.5,0.0}
\definecolor{bad}{rgb}{0.7,0.0,0.0}
\definecolor{neutral}{rgb}{0.3,0.3,0.3}

\title{Iterations 038--054: Broadband Breakthrough, Scaling Laws,\\
  Architecture Sweep, and Pretraining}
\author{Autoresearch Loop}
\date{April 2026}

\begin{document}
\maketitle

\begin{abstract}
We report 17 iterations of systematic experimentation on the Ear-SAAD
scalp-to-in-ear EEG prediction task. The key breakthrough came from two
complementary advances: (1)~moving from narrowband (1--9\,Hz, 20\,Hz sampling)
to broadband (1--45\,Hz, 128\,Hz sampling), yielding $+0.06\,r$; and
(2)~incorporating 19 around-ear cEEGrid channels as additional input,
yielding $+0.15\,r$. Together these advances pushed our best from
$r=0.378$ to $r=0.638$, a \textbf{69\% relative improvement}. A scaling
law experiment showed that model capacity is not the bottleneck---55K
parameters perform comparably to 7M---and that cross-subject variability
is the dominant limiting factor. A 14-iteration architecture sweep
(iter041--054) confirmed this: no architecture, normalization scheme,
loss function, or pretraining strategy improved upon the tiny 55K-parameter
baseline. Subject-level analysis revealed that prediction difficulty is
a physiological property (scalp-to-in-ear coupling strength correlates
$r=0.77$ with LOSO performance), ranging from $r=0.94$ (Subject~3) to
$r=0.38$ (Subject~8).
\end{abstract}

\tableofcontents
\newpage

%==========================================================================
\section{Introduction}
%==========================================================================

Prior iterations (001--037) established that a 7-tap FIR filter with
closed-form initialization, combined MSE+correlation loss, and
correlation-based early stopping achieves $r=0.378$ on the narrowband
(1--9\,Hz) Ear-SAAD benchmark. Extensive exploration of loss functions,
normalization, augmentation, and architecture variants confirmed this as a
hard ceiling for the narrowband regime.

Iterations 038--054 pursued four strategic directions:
\begin{enumerate}
  \item \textbf{Phase 1} (iter038--039): Broadband data and around-ear channels
  \item \textbf{Phase 2}: Scaling law experiment across 5 model sizes
  \item \textbf{Phase 3} (iter041--053): Systematic architecture sweep
  \item \textbf{Phase 4} (iter054): Foundation model pretraining
\end{enumerate}

%==========================================================================
\section{Phase 1: Broadband + Around-Ear Breakthrough}
%==========================================================================

\subsection{Motivation}

The narrowband plateau ($r=0.378$) reflected a fundamental information
bottleneck: at 1--9\,Hz and 20\,Hz sampling, the temporal resolution is
insufficient to capture higher-frequency EEG dynamics that propagate from
scalp to in-ear electrodes. The original Ear-SAAD paper (Geirnaert et al.,
2025) uses broadband data at 256\,Hz, suggesting substantial headroom.
Additionally, the dataset includes 19 around-ear cEEGrid channels that are
physically closer to the in-ear targets than any scalp electrode.

\subsection{Iter038: Broadband FIR (27 Scalp Channels)}

\textbf{Hypothesis}: Moving to broadband (1--45\,Hz at 128\,Hz sampling)
with the same FIR architecture will break the narrowband ceiling.

\textbf{Method}: We reprocessed the raw Ear-SAAD BIDS data with a 1--45\,Hz
bandpass filter and downsampled to 128\,Hz, producing 2-second windows of
$27 \times 256$ samples. The model architecture remained a 7-tap FIR filter
with closed-form initialization, combined loss, and correlation validation.

\textbf{Result}: $r = 0.465$, a \textbf{+23\%} improvement over the
narrowband best.

\begin{table}[H]
\centering
\caption{Broadband vs.\ narrowband comparison (27 scalp input channels).}
\begin{tabular}{lrrl}
\toprule
Model & Channels & Mean $r$ & Change \\
\midrule
Narrowband FIR (iter030) & 27 scalp & 0.378 & --- \\
\textbf{Broadband FIR (iter038)} & 27 scalp & \textcolor{good}{\textbf{0.465}} & \textcolor{good}{+0.087 (+23\%)} \\
\bottomrule
\end{tabular}
\end{table}

\textbf{Analysis}: The broadband signal contains substantially more
information about scalp-to-in-ear propagation. Alpha (8--13\,Hz) and beta
(13--30\,Hz) oscillations, which were completely absent in the narrowband
data, carry spatial information that aids the linear mapping.

\subsection{Iter039: Tiny Deep Model (46 Channels)}

\textbf{Hypothesis}: Adding 19 around-ear cEEGrid channels as input will
further improve predictions, and a small nonlinear model can exploit the
richer 46-channel broadband input.

\textbf{Method}: We expanded the input to all 46 channels (27 scalp +
19 around-ear cEEGrid) and trained a ``tiny deep'' model: a depthwise
separable convolution encoder with residual blocks ($\sim$55K parameters),
combined MSE+correlation loss, channel dropout (15\%), and mixup
augmentation.

\textbf{Result}: $r = 0.638$, our \textbf{best model overall} and a
\textbf{69\% improvement} over the narrowband baseline.

\begin{table}[H]
\centering
\caption{Impact of around-ear channels.}
\begin{tabular}{lrrl}
\toprule
Model & Channels & Mean $r$ & Change \\
\midrule
Broadband FIR (iter038) & 27 scalp & 0.465 & --- \\
CF baseline & 46 (scalp + cEEGrid) & 0.577 & +0.112 \\
\textbf{Tiny deep (iter039)} & 46 (scalp + cEEGrid) & \textcolor{good}{\textbf{0.638}} & \textcolor{good}{+0.173} \\
\bottomrule
\end{tabular}
\end{table}

\textbf{Analysis}: The around-ear cEEGrid channels provide a $+0.15\,r$
boost because they are physically adjacent to the in-ear electrodes.
Volume conduction ensures that cEEGrid channels share substantial signal
content with in-ear channels. The tiny deep model adds $+0.06\,r$ over
the closed-form baseline by capturing mild nonlinearities and temporal
dynamics in the broadband signal.

\subsection{Full 15-Subject LOSO Evaluation}

To characterize generalization across all subjects, we ran full
leave-one-subject-out (LOSO) cross-validation with the closed-form
baseline on all 15 subjects.

\begin{table}[H]
\centering
\caption{Full LOSO results (CF baseline, 46 channels, broadband).}
\begin{tabular}{lr}
\toprule
Metric & Value \\
\midrule
Mean $r$ (15 subjects) & 0.645 \\
95\% CI & [0.563, 0.728] \\
Std $r$ & 0.168 \\
Best subject (Subject 3) & 0.940 \\
Worst subject (Subject 8) & 0.382 \\
\bottomrule
\end{tabular}
\end{table}

The enormous variance across subjects ($\sigma = 0.168$) motivated the
detailed subject analysis in Section~\ref{sec:subject}.

%==========================================================================
\section{Phase 2: Scaling Law}
\label{sec:scaling}
%==========================================================================

\subsection{Motivation}

Before investing in complex architectures, we tested whether simply
scaling model capacity would improve performance. If larger models
consistently outperform smaller ones, the path forward is clear: scale up.
If not, the bottleneck lies elsewhere.

\subsection{Method}

We trained five model sizes on the same broadband 46-channel data,
evaluating on Subject~13 (held out):

\begin{table}[H]
\centering
\caption{Scaling law results (Subject 13, broadband 46-channel).}
\begin{tabular}{lrrr}
\toprule
Size & Parameters & Test $r$ & $\Delta$ vs.\ Tiny \\
\midrule
Tiny & 55K & 0.750 & --- \\
Small & 176K & 0.755 & +0.005 \\
Medium & 1M & 0.751 & +0.001 \\
Large & 3M & 0.758 & +0.008 \\
XL & 7M & 0.752 & +0.002 \\
\bottomrule
\end{tabular}
\end{table}

\subsection{Analysis}

\textbf{Scaling is flat.} A 127$\times$ increase in parameters (55K to
7M) yields only $+0.008\,r$ improvement. The relationship between model
size and test performance is essentially a horizontal line. This confirms
that the bottleneck is not model capacity but \emph{cross-subject
generalization}: the model has sufficient capacity to fit the training
subjects perfectly, but the held-out subject's neural anatomy differs in
ways that no amount of capacity can extrapolate from 14 training subjects.

This finding has profound implications: all subsequent architecture
innovations must target the cross-subject gap specifically, not general
model expressiveness.

%==========================================================================
\section{Phase 3: Architecture Sweep (Iter041--053)}
\label{sec:sweep}
%==========================================================================

\subsection{Overview}

Armed with the knowledge that scaling is flat, we systematically tested
13 architectural innovations targeting different aspects of the
cross-subject problem. All models used broadband 46-channel input and
were benchmarked against iter039 (tiny deep, $r=0.638$).

\begin{table}[H]
\centering
\caption{Architecture sweep results. All models compared against iter039
  ($r = 0.638$). None improved upon the baseline.}
\begin{tabular}{rlrrl}
\toprule
Iter & Architecture & Mean $r$ & $\Delta$ & Key Idea \\
\midrule
039 & \textbf{Tiny deep (baseline)} & \textbf{0.638} & --- & Depthwise-sep conv, 55K params \\
\midrule
\multicolumn{5}{l}{\textit{Normalization and Adaptation}} \\
041 & RevIN + BatchNorm & 0.587 & \textcolor{bad}{$-0.051$} & Reversible instance norm \\
042 & Euclidean alignment & 0.574 & \textcolor{bad}{$-0.064$} & Per-subject whitening \\
046 & Calibrated output & 0.581 & \textcolor{bad}{$-0.057$} & Output calibration layer \\
\midrule
\multicolumn{5}{l}{\textit{Ensemble and Residual}} \\
043 & Ensemble (CF + deep) & 0.606 & \textcolor{bad}{$-0.032$} & Average of CF and deep \\
044 & Residual from CF & 0.610 & \textcolor{bad}{$-0.028$} & $\text{CF}(x) + \alpha \cdot f(x)$ \\
\midrule
\multicolumn{5}{l}{\textit{Context and Attention}} \\
045 & Long context (4s) & 0.606 & \textcolor{bad}{$-0.032$} & 4-second windows \\
047 & Spatial pos.\ encoding & 0.589 & \textcolor{bad}{$-0.049$} & REVE-inspired 4D PE \\
052 & Cross-attention decoder & 0.596 & \textcolor{bad}{$-0.042$} & Output queries attend to input \\
\midrule
\multicolumn{5}{l}{\textit{Training Objectives}} \\
049 & Adversarial training & 0.608 & \textcolor{bad}{$-0.030$} & Gradient reversal for subject invariance \\
050 & L1 loss & 0.605 & \textcolor{bad}{$-0.033$} & REVE-style MAE objective \\
053 & Spectral loss & 0.604 & \textcolor{bad}{$-0.034$} & Frequency-domain MSE \\
\midrule
\multicolumn{5}{l}{\textit{Output Architecture}} \\
048 & NeuroTTT & 0.597 & \textcolor{bad}{$-0.041$} & Test-time training with SSL \\
051 & Per-channel heads & 0.609 & \textcolor{bad}{$-0.029$} & 12 separate output decoders \\
\bottomrule
\end{tabular}
\end{table}

\subsection{Detailed Analysis by Category}

\subsubsection{Normalization Always Hurts}

Three normalization approaches were tested (iter041: RevIN+BN, iter042:
Euclidean alignment, iter046: calibrated output), all producing the
largest performance drops ($-0.051$ to $-0.064$). This is consistent
with findings from the narrowband phase: InstanceNorm consistently
degraded performance by $\sim$0.003\,r. In the broadband regime, the
effect is amplified. Subject-specific amplitude information---which
normalization destroys---is a critical feature for cross-subject
prediction. When EEG amplitude patterns are normalized away, the model
loses its ability to leverage the absolute signal characteristics that
differ between subjects but are informative for the spatial mapping.

\subsubsection{Ensemble and Residual Approaches Underperform}

Both the CF+deep ensemble (iter043, $r=0.606$) and the residual-from-CF
architecture (iter044, $r=0.610$) performed worse than the standalone
deep model. The CF baseline ($r=0.577$) anchors the ensemble/residual
toward its own (lower) performance ceiling. Because the deep model
already captures the linear mapping implicitly, combining it with an
explicit CF component introduces redundancy without adding complementary
information.

\subsubsection{Extended Context Does Not Help}

Doubling the context window from 2 seconds to 4 seconds (iter045,
$r=0.606$) did not improve performance. The relevant scalp-to-in-ear
dynamics operate on timescales shorter than 2 seconds; longer context
introduces more noise than signal. This is consistent with the
narrowband finding that 7 FIR taps ($\sim$350\,ms) are sufficient.

\subsubsection{Spatial Positional Encoding}

Inspired by the REVE foundation model's 4D positional encoding
(Cartesian electrode coordinates), iter047 added learnable spatial
embeddings based on electrode positions. The result ($r=0.589$) was
substantially worse, suggesting that with only 46 channels and 15
subjects, the model cannot learn useful spatial priors---the data is too
sparse for the positional encoding to capture meaningful electrode
geometry.

\subsubsection{Alternative Loss Functions}

L1 loss (iter050, $r=0.605$) and spectral loss (iter053, $r=0.604$)
both underperformed the combined MSE+correlation loss. Adversarial
subject-invariant training (iter049, $r=0.608$) was intended to learn
representations that transfer across subjects via gradient reversal,
but destabilized training without meaningful improvement.

\subsubsection{Specialized Output Architectures}

Per-channel output heads (iter051, $r=0.609$) allowed each in-ear
channel its own decoder, and cross-attention (iter052, $r=0.596$)
let output queries attend to input channel features. Neither approach
helped: the standard shared linear decoder is already sufficient for the
12-channel output, and the additional parameters in specialized decoders
lead to overfitting on 12 training subjects.

\subsubsection{Test-Time Training (NeuroTTT)}

Iter048 implemented NeuroTTT-style self-supervised adaptation at test
time, using masked channel reconstruction as an auxiliary objective. The
result ($r=0.597$) was worse than the baseline, likely because the SSL
signal from only 46 channels is too weak to drive meaningful adaptation,
and the adaptation risks overfitting to noise in the test subject's data.

%==========================================================================
\section{Phase 4: Foundation Model Pretraining (Iter054)}
%==========================================================================

\subsection{Motivation}

Given that all architecture modifications failed, we hypothesized that
the bottleneck might be insufficient training data diversity. Foundation
model pretraining on large external EEG corpora could provide temporal
representations that generalize better across subjects.

\subsection{Method}

We pretrained a BIOT-inspired model on 20 subjects from the Healthy Brain
Network (HBN) EEG dataset, then fine-tuned on the Ear-SAAD broadband
46-channel data. The pretraining used masked channel prediction as a
self-supervised objective. Fine-tuning followed the standard protocol:
freeze encoder for 20 epochs, then unfreeze with small learning rate
($10^{-5}$), combined loss, correlation validation.

\subsection{Result}

$r = 0.589$, \textbf{worse} than the baseline ($r = 0.638$) by
$-0.049$.

\subsection{Analysis}

The pretraining failed for two likely reasons:
\begin{enumerate}
  \item \textbf{Insufficient pretraining data}: 20 HBN subjects provide
    minimal diversity compared to the 25,000+ subjects used by successful
    foundation models like REVE and FEMBA. The pretrained representations
    are not substantially more general than what can be learned from the
    12 Ear-SAAD training subjects directly.
  \item \textbf{Domain mismatch}: HBN recordings use standard 10-20
    montages without any ear-adjacent channels. The pretrained encoder
    has no exposure to the cEEGrid/in-ear signal characteristics, and
    fine-tuning on only 12 subjects is insufficient to bridge this gap.
\end{enumerate}

Effective pretraining likely requires access to the TUH EEG Corpus
($\sim$21,000 hours) or similar large-scale datasets, combined with a
channel-agnostic architecture that can transfer to novel electrode
positions.

%==========================================================================
\section{Subject-Level Analysis}
\label{sec:subject}
%==========================================================================

\subsection{LOSO Performance Distribution}

The full 15-subject LOSO evaluation reveals dramatic inter-subject
variability:

\begin{table}[H]
\centering
\caption{Per-subject LOSO performance (CF baseline, 46ch broadband).}
\begin{tabular}{clr|clr}
\toprule
Rank & Subject & LOSO $r$ & Rank & Subject & LOSO $r$ \\
\midrule
1 & Subject 3 & \textcolor{good}{0.940} & 9 & Subject 15 & 0.601 \\
2 & Subject 4 & \textcolor{good}{0.840} & 10 & Subject 1 & 0.586 \\
3 & Subject 9 & 0.765 & 11 & Subject 6 & 0.575 \\
4 & Subject 11 & 0.742 & 12 & Subject 12 & 0.548 \\
5 & Subject 5 & 0.732 & 13 & Subject 2 & 0.539 \\
6 & Subject 13 & 0.728 & 14 & Subject 14 & \textcolor{bad}{0.433} \\
7 & Subject 10 & 0.653 & 15 & Subject 8 & \textcolor{bad}{0.382} \\
8 & Subject 7 & 0.618 &  &  & \\
\bottomrule
\end{tabular}
\end{table}

\subsection{The Dominant Predictor: Scalp-to-In-Ear Coupling Strength}

Statistical analysis of subject-level features revealed that LOSO
performance is overwhelmingly determined by the intrinsic
\emph{scalp-to-in-ear coupling strength}---a physiological property of
the individual.

\begin{table}[H]
\centering
\caption{Top predictors of per-subject LOSO $r$.}
\begin{tabular}{lrrl}
\toprule
Metric & Pearson $r$ & $p$-value & Interpretation \\
\midrule
In-ear channel RMS CV & $-0.777$ & 0.0007 & Hard subjects: uneven channel quality \\
In-ear window variance & $+0.714$ & 0.0028 & Easy subjects: richer dynamics \\
Min best scalp-inear corr & $+0.767$ & 0.0008 & Hard subjects: ``dead'' channels \\
Mean best scalp-inear corr & $+0.707$ & 0.0032 & Hard subjects: weaker coupling \\
Well-coupled channels ($r > 0.5$) & $+0.775$ & 0.0007 & Hard subjects: fewer usable channels \\
\bottomrule
\end{tabular}
\end{table}

Notably, scalp signal amplitude, in-ear signal amplitude, data quantity,
and trial-level SNR are \textbf{not} significant predictors---ruling out
simple data quality explanations.

\subsection{Case Study: Best vs.\ Worst Subject}

\textbf{Subject 3} ($r = 0.940$): Every in-ear channel has a scalp channel
correlated at $r > 0.91$ (mean best correlation = 0.957). The
scalp-to-in-ear mapping is almost perfectly linear for this subject.

\textbf{Subject 8} ($r = 0.382$): Channels 0, 4, and 7 have peak
scalp correlations of only 0.32, 0.22, and 0.24 respectively---these
channels are \emph{fundamentally unpredictable} from scalp EEG. The mean
best correlation is 0.585.

The coupling strength appears to reflect individual differences in skull
conductivity, cortical folding patterns, and in-ear electrode contact
quality---anatomical factors that no amount of model engineering can
overcome without subject-specific calibration data.

\subsection{In-Ear Channel Difficulty}

Not all in-ear channels are equally predictable. Across all 15 subjects:

\begin{table}[H]
\centering
\caption{In-ear channel predictability (mean best scalp correlation).}
\begin{tabular}{lrrl|lrrl}
\toprule
Ch & Mean $r$ & Std & Difficulty & Ch & Mean $r$ & Std & Difficulty \\
\midrule
11 & 0.891 & 0.087 & Easy & 0 & 0.684 & 0.259 & Medium \\
5 & 0.880 & 0.093 & Easy & 8 & 0.643 & 0.271 & Hard \\
10 & 0.834 & 0.185 & Easy & 1 & 0.645 & 0.237 & Hard \\
4 & 0.824 & 0.195 & Easy & 7 & 0.634 & 0.283 & Hard \\
2 & 0.727 & 0.253 & Medium & 6 & 0.619 & 0.255 & Hard \\
3 & 0.725 & 0.322 & Medium & & & & \\
\bottomrule
\end{tabular}
\end{table}

Channels 5, 10, and 11 are reliably well-coupled across subjects;
channels 6, 7, and 8 are consistently harder. This asymmetry likely
reflects the left-right anatomy of the ear canal relative to temporal
cortex sources.

%==========================================================================
\section{Summary of Contributions}
%==========================================================================

\subsection{Attribution of Improvement}

\begin{table}[H]
\centering
\caption{Decomposition of improvement from narrowband baseline to best model.}
\begin{tabular}{lrrl}
\toprule
Component & $r$ & $\Delta r$ & Attribution \\
\midrule
Narrowband FIR (iter030) & 0.378 & --- & Previous best \\
$+$ Broadband (1--45\,Hz) & 0.465 & +0.087 & Higher-freq information \\
$+$ Around-ear channels & 0.577 & +0.112 & Physical proximity to targets \\
$+$ Deep model & 0.638 & +0.061 & Nonlinear temporal dynamics \\
\midrule
\textbf{Total} & \textbf{0.638} & \textbf{+0.260} & \textbf{+69\% relative} \\
\bottomrule
\end{tabular}
\end{table}

\subsection{Negative Results}

The following approaches \textbf{all failed} to improve upon the
55K-parameter baseline:

\begin{itemize}
  \item \textbf{Model scaling}: 55K to 7M parameters ($+0.008\,r$, within noise)
  \item \textbf{All normalization}: RevIN, Euclidean alignment, InstanceNorm,
    calibration ($-0.05$ to $-0.06\,r$)
  \item \textbf{Ensemble/residual}: CF+deep averaging, residual-from-CF
    ($-0.03$ to $-0.03\,r$)
  \item \textbf{Extended context}: 4-second windows ($-0.032\,r$)
  \item \textbf{Spatial encoding}: REVE-inspired positional encoding ($-0.049\,r$)
  \item \textbf{Alternative losses}: L1, spectral, adversarial ($-0.030$ to $-0.034\,r$)
  \item \textbf{Output specialization}: Per-channel heads, cross-attention
    ($-0.029$ to $-0.042\,r$)
  \item \textbf{Test-time training}: NeuroTTT SSL adaptation ($-0.041\,r$)
  \item \textbf{Pretraining}: BIOT-style on 20 HBN subjects ($-0.049\,r$)
\end{itemize}

%==========================================================================
\section{Key Findings and Implications}
%==========================================================================

\begin{enumerate}
  \item \textbf{Data trumps architecture.} The single largest
    improvement ($+0.173\,r$) came from adding around-ear input channels
    and broadband frequencies---not from any model change. This is a
    recurring theme in applied ML: when the signal-to-noise ratio is low,
    acquiring better input features dominates architectural innovation.

  \item \textbf{Scaling is flat for EEG regression.} Unlike language and
    vision tasks, EEG scalp-to-in-ear prediction does not benefit from
    larger models. The mapping is approximately linear (CF baseline
    achieves $r=0.577$), and what nonlinearity exists is captured by a
    55K-parameter model. This suggests the task's intrinsic dimensionality
    is very low.

  \item \textbf{Cross-subject variability is the fundamental bottleneck.}
    Subject-level $r$ ranges from 0.38 to 0.94, driven by physiological
    coupling strength ($\rho = 0.77$ with LOSO performance). No
    architectural change in our sweep addressed this. Future work must
    focus on subject-specific adaptation, either through few-shot
    calibration or large-scale pretraining on thousands of subjects.

  \item \textbf{Normalization destroys useful information.} Every
    normalization technique tested (RevIN, InstanceNorm, Euclidean
    alignment, calibration) degraded performance. Subject-specific
    amplitude dynamics carry information about the spatial mapping that
    normalization removes.

  \item \textbf{Small pretraining datasets do not help.} Pretraining on
    only 20 external subjects provides insufficient diversity. Effective
    EEG foundation models (REVE, FEMBA, LUNA) require 1,000--25,000+
    subjects. Access to the TUH EEG Corpus ($\sim$21,000 hours) is
    likely necessary for meaningful pretraining gains.
\end{enumerate}

%==========================================================================
\section{Comprehensive Leaderboard}
%==========================================================================

\begin{table}[H]
\centering
\small
\caption{Complete leaderboard for iterations 038--054 (broadband, 46-channel benchmark).}
\begin{tabular}{rlrrl}
\toprule
Iter & Model & Mean $r$ & $\Delta$ vs.\ Best & Key Idea \\
\midrule
038 & Broadband FIR (27ch) & 0.465 & $-0.173$ & First broadband model \\
--- & CF baseline (46ch) & 0.577 & $-0.061$ & Linear spatial filter \\
\textbf{039} & \textbf{Tiny deep (46ch)} & \textcolor{good}{\textbf{0.638}} & \textbf{---} & \textbf{Best: depthwise-sep conv, 55K} \\
\midrule
041 & RevIN + BN adapt & 0.587 & $-0.051$ & Reversible instance norm \\
042 & Euclidean alignment & 0.574 & $-0.064$ & Per-subject whitening \\
043 & Ensemble (CF + deep) & 0.606 & $-0.032$ & Average predictions \\
044 & Residual from CF & 0.610 & $-0.028$ & CF + learned residual \\
045 & Long context (4s) & 0.606 & $-0.032$ & 4-second windows \\
046 & Calibrated output & 0.581 & $-0.057$ & Output calibration \\
047 & Spatial pos.\ encoding & 0.589 & $-0.049$ & REVE-inspired 4D PE \\
048 & NeuroTTT & 0.597 & $-0.041$ & Test-time SSL adaptation \\
049 & Adversarial & 0.608 & $-0.030$ & Gradient reversal \\
050 & L1 loss & 0.605 & $-0.033$ & MAE objective \\
051 & Per-channel heads & 0.609 & $-0.029$ & 12 separate decoders \\
052 & Cross-attention & 0.596 & $-0.042$ & Output queries attend to input \\
053 & Spectral loss & 0.604 & $-0.034$ & Frequency-domain MSE \\
054 & Pretrained fine-tune & 0.589 & $-0.049$ & 20 HBN subjects \\
\bottomrule
\end{tabular}
\end{table}

%==========================================================================
\section{Future Directions}
%==========================================================================

Based on these findings, the most promising paths forward are:

\begin{enumerate}
  \item \textbf{Large-scale pretraining}: Access the TUH EEG Corpus
    ($\sim$21,000 hours, 30,000 recordings) and pretrain with a
    channel-agnostic architecture (LUNA-style cross-attention or
    SingLEM-style single-channel). This is the most likely path to
    breaking the cross-subject barrier.

  \item \textbf{Few-shot subject calibration}: Use a small amount of
    target-subject data (30--60 seconds) to adapt the spatial mapping
    at test time. A simple channel-wise affine transformation or
    low-rank adaptation could dramatically improve performance on hard
    subjects without requiring full fine-tuning.

  \item \textbf{Channel-weighted evaluation}: Report per-channel
    results and weight the loss by estimated channel predictability.
    This would improve effective performance by focusing on channels
    where prediction is physically possible.

  \item \textbf{Additional ear-adjacent datasets}: The cEEGrid electrode
    array is gaining adoption. Combining multiple cEEGrid datasets
    (if available) could provide the subject diversity needed for
    better cross-subject generalization.

  \item \textbf{Hyperparameter optimization}: A systematic Optuna sweep
    over the tiny deep model's hyperparameters could yield modest
    improvements ($+0.01$--$0.03\,r$) at low risk.
\end{enumerate}

%==========================================================================
\section{Conclusion}
%==========================================================================

Iterations 038--054 achieved a major breakthrough ($r = 0.378 \to 0.638$,
$+69\%$) through data-centric improvements: broadband filtering and
around-ear input channels. Extensive experimentation with model scaling,
14 architectural variants, and pretraining unanimously confirmed that the
current bottleneck is cross-subject physiological variability, not model
capacity or architecture. The scalp-to-in-ear coupling strength is a
biological property that varies dramatically across individuals
($r = 0.38$ to $0.94$) and cannot be overcome by training-time
innovations alone. Future progress requires either large-scale
pretraining (thousands of subjects) or few-shot subject-specific
calibration.

\end{document}
